% !PRINTED 436
\weberpage{436}
%<<CONT
one also sees easily that these $n$ values are all different from one another, if one leaves aside the special case that $\varphi$ is a multiple of $\pi$. For two cosines are equal to one another only when the sum or the difference of the angles is a multiple of $2\pi$.

The equation
\weq{15}{\sin n\varphi = \sin\varphi\, B_n(x)}
determines $\sin\varphi$ uniquely in terms of $x$, when $B_n(x)$ is not equal to zero.

If we seek in equation (13) the term independent of $x$, then with the use of (6) there results for an even $n$
\[
2(-1)^{\frac{n}{2}} - 2\cos n\varphi,
\]
and for an odd $n$
\[
-2\cos n\varphi,
\]
and if one sets this equal to the positive or negative product of the roots, one obtains for an even $n$
\weq{16}{2^{n-1} \overset{\nu}{\Pi} \cos\left(\varphi + \dfrac{2\nu\pi}{n}\right) = (-1)^{\frac{n}{2}} - \cos n\varphi,}
and for an odd $n$
\weq{17}{2^{n-1} \overset{\nu}{\Pi} \cos\left(\varphi + \dfrac{2\nu\pi}{n}\right) = \cos n\varphi,}
and these formulae hold, since on the right and on the left stand continuous functions of $\varphi$, also for the case at first excluded, where $\varphi$ is a multiple of $\pi$.

Herein we may put $\nu m$ for $\nu$ under the product sign, if $m$ is any number coprime to $n$, because then $\nu m$ runs, together with $\nu$, through a complete system of residues to the modulus $n$. If we separate off the value $\nu=0$, and then let $\nu$ run through both the positive and the negative numbers that are absolutely smaller than $n:2$, we obtain from (17) the formula, valid for every odd $n$,
\weq{18}{2^{n-1} \overset{\nu}{\underset{1,\,\frac{n-1}{2}}{\Pi}} \cos\left(\varphi + \dfrac{2\nu m\pi}{n}\right) \cos\left(\varphi - \dfrac{2\nu m\pi}{n}\right) = \dfrac{\cos n\varphi}{\cos\varphi}.}
If we set $\varphi = 0$ herein and take the square root, there results
%CONT>>
